\chapter{Barium Swallow}

\section*{What Is This Test?}
A barium swallow is an imaging examination of swallowing and the upper gastrointestinal tract. The patient drinks a contrast liquid called barium while a series of X-ray images, usually fluoroscopy, records its movement through the mouth, throat, and esophagus.

\section*{Alternative Names}
\begin{itemize}
  \item Esophagram
  \item Upper Gastrointestinal Contrast Study
  \item Modified Barium Swallow (when focused on the oropharyngeal phase)
\end{itemize}

\section*{Principle}
Barium coats the lining of the digestive tract and blocks X-rays, outlining anatomy and showing transit. A modified barium swallow evaluates the oral and pharyngeal phases with a speech-language pathologist and radiologist.

\section*{Why This Test Is Ordered}
It may evaluate difficulty swallowing, food sticking, aspiration, reflux-related symptoms, suspected narrowing, rings, diverticula, ulcers, or abnormal esophageal movement.

\section*{Primary vs Secondary Test}
The study is selected after clinical assessment. Endoscopy may be preferred when mucosal inspection, biopsy, dilation, or treatment is likely to be needed; a swallow study is particularly useful for dynamic swallowing and aspiration assessment.

\section*{How This Test Is Performed}
The patient drinks barium of varying consistencies while images are taken in multiple positions. The study commonly takes 15--30 minutes. Barium is not absorbed and normally passes through the stool.

\section*{What Preparation Is Needed}
Fasting is commonly required for several hours, according to local instructions. Tell the imaging team about pregnancy possibility, severe constipation, bowel obstruction, or a prior reaction to contrast materials.

\section*{Contraindications and Precautions}
There is a small radiation exposure. Aspiration can occur in patients with severe swallowing impairment. Barium may cause constipation and should not be used when gastrointestinal perforation is strongly suspected; a water-soluble contrast agent may be selected instead.

\section*{Understanding Results}
The report may describe delayed transit, aspiration, reflux, narrowing, a mass, diverticulum, abnormal motility, or structural compression. Results must be correlated with symptoms because a normal study does not exclude every esophageal or gastric disorder.

\section*{Follow-up Tests}
Follow-up may include upper endoscopy, esophageal manometry, pH monitoring, CT or MRI, speech and swallowing therapy, or surgical consultation depending on the finding.

\section*{References}
\begin{enumerate}
  \item MedlinePlus. Barium Swallow.
  \item American College of Radiology. Practice parameter for the performance of contrast esophagrams.
\end{enumerate}

\clearpage
\section*{Understanding Results and Follow-up}
The laboratory report should identify the specimen, method, units, reference interval or decision limit, and whether the result is positive, negative, abnormal, or inconclusive. Reference intervals vary with age, sex, pregnancy, specimen type, assay, and laboratory; a result outside a range is not automatically a diagnosis, and a result within a range does not always exclude disease. Discuss unexpected results with the ordering clinician, who can decide whether repeat or confirmatory testing, treatment, or referral is appropriate. Do not change medicines or delay urgent care based on an isolated result.


\section*{Regulatory Status and Authorization Date}
This chapter describes a test category, not a specific commercial product. FDA, CDSCO, EU IVDR, and MHRA approval or registration dates therefore cannot be assigned without the assay manufacturer, product name, instrument, and intended use. For laboratory-developed tests, an individual agency approval date may not exist. See the Preface for the interpretation of regulatory status.

\clearpage
